\clearpage
\section{Complete affine-lattice statement from the short note}
\label{long251:app:affine-detail}
The following expands the compact circularity statement formerly printed
in the body of this record.  In particular the fixed-lattice predicate
and its strict growth hypothesis are stated, not left implicit.

Here $D_N=\sigma_h(N)$ and
$\delta_N=g_{N+h+1}-g_{N+1}$, with $h$ fixed.  For a longer block, set
\[
 B_{h,N,r}=\sum_{i=0}^{r-1}2^{r-1-i}\delta_{N+i};
 \qquad D_{N+r}=2^rD_N-B_{h,N,r}.
\]
Call the data-dependent affine condition
\[
 \mathcal A_{N,r}:\qquad
 D_{N+r}\in -B_{h,N,r}+2^{r+1}\Z.
\]

\begin{theorem}[affine and fixed-lattice circularity]\label{long251:res:affinecollapse}\compatlabel{res:affinecollapse}
For every rational dyadic tail recurrence and all $h,N,r\ge0$,
\begin{equation}
 \mathcal A_{N,r}\quad\Longleftrightarrow\quad D_N\in2\Z.
\label{eq:affinecollapse}\compatlabel{long251:eq:affinecollapse}
\end{equation}
Consequently, if every $\delta_N$ is even, then
\begin{equation}
 \bigl(\forall N_0\ \exists N,r:\ N_0<N\text{ and }\neg\mathcal A_{N,r}\bigr)
 \quad\Longleftrightarrow\quad
 D_N\notin\Z\text{ for arbitrarily large }N.
\label{eq:affinecofinal}\compatlabel{long251:eq:affinecofinal}
\end{equation}

There is a second equivalence.  Let $b:\N\to\Q$ satisfy
$|D_N|\le b(N)$ for every $N$, and suppose that for every $N$ and every
positive integer $q$ there is an $r$ with
\begin{equation}
 2b(N+r)q<2^r.
\label{eq:dyadicscale}\compatlabel{long251:eq:dyadicscale}
\end{equation}
Then
\begin{equation}
 \begin{split}
 &\forall N_0\ \exists N,r:\ N_0<N\text{ and }
   \forall z\in\Z,\quad
   b(N+r)<|B_{h,N,r}-2^rz|\\
 &\hspace{35mm}\Longleftrightarrow\quad
 D_N\notin\Z\text{ for arbitrarily large }N.
 \end{split}
\label{eq:fixedcollapse}\compatlabel{long251:eq:fixedcollapse}
\end{equation}
\end{theorem}

\begin{proof}
Substituting $D_{N+r}=2^rD_N-B_{h,N,r}$ into $\mathcal A_{N,r}$ cancels the
observed block from both sides and leaves $D_N=2z$.  This proves
\eqref{long251:eq:affinecollapse}.  After one recurrence step, evenness of
$\delta_N$ identifies even integrality at $N+1$ with ordinary integrality at
$N$.  The forward implication in~\eqref{long251:eq:affinecofinal} follows, while the
reverse implication chooses a nonintegral $D_N$ and the legal depth $r=0$.

For~\eqref{long251:eq:fixedcollapse}, eventual integrality and the block identity give
some $z\in\Z$ with
\[
 B_{h,N,r}-2^rz=-D_{N+r},
\]
so the displayed strict separation contradicts $|D_{N+r}|\le b(N+r)$.
Conversely, if $D_N$ is nonintegral and has reduced denominator $q$, then its
distance from every integer is at least $1/q$.  Choose $r$ from
\eqref{long251:eq:dyadicscale}, scale this separation by $2^r$, and use the block
identity together with $|D_{N+r}|\le b(N+r)$.  The triangle inequality gives
$|B_{h,N,r}-2^rz|>b(N+r)$ for every integer $z$.
\end{proof}

\noindent The three displayed equivalences are assembled in one declaration,
\href{https://github.com/wcook04/plectis-erdos/blob/3d6d938d696fed0fb71dd55115a18a73738ff223/lean/ErdosProblems/Erdos251/PaperCoreR7.lean\#L241}{affine circularity bundle}.

The first equivalence shows that the apparent affine hierarchy contains no
depth-dependent information: it is the pullback of even integrality through
the recurrence identity.  The fixed lattice removes that data dependence, but
under the stated growth condition rational denominator separation already
forces every required escape.  Hence neither cofinal condition is an
independent source of information about consecutive primes.

\section{Mathematical and evidence map}\label{sec:erdos-251-complete-family-map}
\paragraph{Perturbations.}
The bounded construction and the sparse congruence-preserving theorem are
existence results for altered integer words.  Their conclusion is not a
prime-producing construction.  The ordinary sparse proof and the formal
existence proof have different witnesses.
\paragraph{Exact criteria.}
Summation by parts identifies the actual tails.  Integral-shift, free-pair
and lcm-diagonal criteria then classify rationality at the recurrence level.
Their equivalence does not supply the missing prime-specific witnesses.
\paragraph{Obstructions.}
Polynomial growth, bounded residues, fixed-block nonconcentration and
recurring gap values do not by themselves force the needed tail behaviour.
The numbered constructions above give the precise counterexamples.
\paragraph{Certificates.}
Finite window comparisons and denominator exclusions are exact finite
statements once the proved analytic remainder bounds are supplied.  A
large rational-denominator floor is not an irrationality theorem.
\paragraph{Formal evidence.}
The revision's declaration ledger records the file, line, immutable
revision and supplied build status of each linked declaration.  The older
administrative family catalogue is retained verbatim with that ledger;
counts of registry rows or project-wide families do not strengthen any
mathematical statement.

